\documentclass{article}
\usepackage{tcee-repaso}
\tcetema[Procesos de integración no comunitarios]{3.B.35}{Procesos de integración no comunitarios.}
\tcefecha{19/11/2026}
\begin{document}
\cabecerarepaso
\begin{columnas}

% ==========================================================================
\bloque[2,5]{Introducción}
\begin{itemize}
\item \textbf{Enganche}: «más vale solo que mal acompañado»: integrarse o no, y con quién.
  \begin{itemize}
  \item Fuera de la UE, a medias: se quitan aranceles, pero conservan política comercial y soberanía.
  \item 2026, las 2 respuestas: UE-Mercosur provisional (1-5) y ZLECAf con 50 ratificaciones; EE.~UU. no renueva el T-MEC (1-7: revisiones anuales).
  \item ¿Abren los bloques el comercio mundial o lo trocean?
  \end{itemize}
\item \textbf{Relevancia}:
  \begin{itemize}
  \item \dato{390 ACR en vigor}{OMC, 2026}; todo miembro de la OMC tiene alguno.
  \item T-MEC $\approx$30\,\% del PIB mundial; RCEP 27\,\%.
  \item España negocia por la UE \vertema{3.B.41}; América Latina \vertema{4.A.23}.
  \end{itemize}
\item \textbf{Contextualización}: \concepto{integración}: proceso y estado sin discriminación (\autor{Balassa}, 1961).
  \begin{itemize}
  \item 4 olas: \emph{cerrado} (50-70); \emph{abierto} (80-90); bilaterales y profundos (1996-2016); \emph{megaacuerdos} y fragmentación (2017-).
  \item Teoría \vertema{3.B.10}; UE \vertema{3.B.36}; sus acuerdos \vertema{3.B.41}.
  \end{itemize}
\item \textbf{Problemática}: ¿por qué y con quién? ¿Peldaño o escollo? ¿Qué procesos hay y cuánto integran?
\item \textbf{Estructura}: \hyperref[rep:I]{I} teoría y reglas; \hyperref[rep:II]{II} los casos, región a región; \hyperref[rep:III]{III} balance.
\end{itemize}
\esquemacante{%
\apartadocante{Introducción}{2,5}
\apartadocante{I. Justificación y marco jurídico}{7}
\subcante{1. Formas de integración}{1,5}
\subcante{2. ¿Por qué se integran?}{1}
\subcante{3. ¿Peldaño o escollo?}{3,5}
\subcante{4. Marco jurídico}{1}
\apartadocante{II. Los procesos no comunitarios}{13}
\subcante{1. Europa y Eurasia}{2}
\subcante{2. América}{4,5}
\subcante{3. Asia-Pacífico}{3,5}
\subcante{4. Oriente Medio y norte de África}{1}
\subcante{5. África}{2}
\apartadocante{III. Balance y situación actual}{6,5}
\subcante{1. Efectos económicos}{3,5}
\subcante{2. Un mundo fragmentado}{3}
\apartadocante{Conclusión}{1}}

% ==========================================================================
\bloque[7]{I. Justificación y marco jurídico\label{rep:I}}
\begin{ideaclave}
La integración regional es un segundo óptimo: libera entre los socios y discrimina a los demás. Si es peldaño o escollo depende de su diseño: abierta y con arancel exterior bajo, peldaño; cerrada, escollo.
\end{ideaclave}
\graficorepaso[\linewidth]{olas-regionalismo}{Las 4 olas del regionalismo, 1957-2026}

\sub{I.1. Formas de integración\label{rep:I.1}}
\begin{itemize}
\item \emph{Escalera de \autor{Balassa}}:
  \begin{itemize}
  \item Preferencias (D-8) $\to$ \concepto{ZLC} (T-MEC, AELC, CPTPP, RCEP, ZLECAf) $\to$ \concepto{UA} (Mercosur, CAN, CCG, SACU, UEE, EAC) $\to$ mercado común $\to$ UEM (UEMOA, CEMAC) $\to$ integración total.
  \end{itemize}
\item \emph{ZLC frente a UA}:
  \begin{itemize}
  \item ZLC: \concepto{desviación del tráfico} (\emph{trade deflection}) $\Rightarrow$ \concepto{reglas de origen}.
  \item UA: no las necesita, pero exige ceder la política comercial.
  \item Por eso casi todo fuera de la UE son ZLC: 390 ACR = 343 ZLC, 26 alcance parcial, 16 UA, 5 servicios.
  \end{itemize}
\item \emph{Integración negativa/positiva} (\autor{Tinbergen}, 1954). \emph{Profundidad}: OMC-plus y OMC-extra (\autor{Horn}-\autor{Mavroidis}-\autor{Sapir}, 2010).
\end{itemize}
\graficorepaso[0.85\linewidth]{zlc-ua}{ZLC: el producto entra por el socio de arancel bajo; UA: una sola frontera}

\sub{I.2. ¿Por qué se integran los países?\label{rep:I.2}}
\begin{itemize}
\item \emph{Primer óptimo}: libre comercio multilateral. \emph{Segundo óptimo}:
  \begin{itemize}
  \item \autor{Viner} (1950): creación frente a desviación \vertema{3.B.10}.
  \item \autor{Cooper}-\autor{Massell} (1965): la rebaja no discriminatoria domina.
  \item \autor{Kemp}-\autor{Wan} (1976): siempre hay un arancel común que no daña a terceros.
  \end{itemize}
\item \emph{Lo que la justifica}:
  \begin{itemize}
  \item RRI (\autor{Mundell}, 1964); escala, competencia, IED.
  \item Economía política: desviar ayuda a aprobar (\autor{Grossman}-\autor{Helpman}, 1995).
  \item No tradicionales: compromiso, señal, seguro, poder negociador, seguridad (\autor{Fernández}-\autor{Portes}, 1998).
  \end{itemize}
\end{itemize}

\sub{I.3. ¿Peldaño o escollo?\label{rep:I.3}}
\begin{itemize}
\item \autor{Summers} (1992): «todos los ismos»; \autor{Bhagwati}: el regionalismo mina el GATT.
  \begin{itemize}
  \item 2 preguntas (\autor{Bhagwati}, 1993): \emph{estática} (\autor{Viner}) y \emph{dinámica}. Liberalización \emph{discriminatoria}.
  \end{itemize}
\end{itemize}
\subsub{a) Argumentos del peldaño}
\begin{itemize}
\item \autor{Kemp}-\autor{Wan}: cada ampliación, un paso al libre comercio (lógica del art.~XXIV).
\item \emph{Socios naturales} (\autor{Summers}; \autor{Krugman}, 1991b): poca desviación; críticas de \autor{Bhagwati}-\autor{Panagariya} (1996).
\item \concepto{Dominó} (\autor{Baldwin}, 1993): quien queda fuera sufre la desviación y presiona para entrar.
  \begin{itemize}
  \item Equilibrio en $P(n^*)=R(n^*)$; explica las olas (Austria, Suecia, Finlandia; Canadá en el TLCAN).
  \item Si el bloque admite a quien lo pida, casi multilateral.
  \end{itemize}
\end{itemize}
\graficorepaso[0.75\linewidth]{domino-baldwin}{Dominó: presión para entrar ($P$) y resistencia ($R$) según el tamaño del bloque}
\begin{itemize}
\item \emph{Apisonadora} (2015): cada rebaja encoge a los que compiten con importaciones $\Rightarrow$ más apoyo a la siguiente.
\item \emph{Complementariedad arancelaria} (\autor{Ornelas}, 2005): ZLC $\Rightarrow$ interés en bajar el arancel exterior (América Latina, \autor{Estevadeordal}-\autor{Freund}-\autor{Ornelas}, 2008).
\item \emph{Laboratorio de reglas} (servicios, inversión, digital) y \emph{multilateralización} (\autor{Baldwin}, 2006): origen acumulable.
\end{itemize}
\subsub{b) Argumentos del escollo}
\begin{itemize}
\item \concepto{Bloques de \autor{Krugman}} (1991): menos bloques = más comercio libre, pero más poder de mercado.
  \begin{itemize}
  \item Arancel de \autor{Nash}: $t^N\approx1/[(\sigma-1)(1-\tilde s)]$.
  \item Bienestar mundial \textbf{mínimo con 3 bloques}; con 1, libre comercio.
  \item Lección: el problema es el arancel exterior, no el regionalismo.
  \end{itemize}
\end{itemize}
\graficorepaso[0.8\linewidth]{krugman-bloques}{Bienestar mundial según el número de bloques: mínimo con 3 para $\sigma=2,4,8$}
\begin{itemize}
\item \concepto{Plato de espaguetis} (\autor{Bhagwati}, 1995, 2008): reglas de origen superpuestas; a veces sale más barato pagar la NMF.
\item \emph{Intereses creados}:
  \begin{itemize}
  \item \autor{Krishna} (1998): los ACR desviadores tienen más apoyo y luego bloquean lo multilateral ($\Delta\Pi^{M|ACR}=-3t^2/8<0$).
  \item Votante mediano (\autor{Levy}, 1997); EE.~UU. rebajó menos donde daba preferencias (\autor{Limão}, 2006).
  \end{itemize}
\item \emph{Centro y radios} (\autor{Kowalczyk}-\autor{Wonnacott}, 1992); distracción y poder del grande.
\end{itemize}
\subsub{c) Balance: cómo hacer del regionalismo un peldaño}
\begin{itemize}
\item Creación $>$ desviación y NMF a la baja, pero Doha atascada.
\item Peldaño si: arancel exterior bajo, puerta abierta, origen acumulable, reglas profundas no discriminatorias, vigilancia multilateral.
\item Propuestas: bajar la NMF (\autor{Bhagwati}), multilateralizar (\autor{Baldwin}), plurilaterales en la OMC.
\end{itemize}

\sub{I.4. Marco jurídico\label{rep:I.4}}
\begin{itemize}
\item \emph{Art.~XXIV GATT}: ZLC y UA si:
  \begin{itemize}
  \item liberan «lo esencial de los intercambios» ($\approx$90\,\%);
  \item en un plazo razonable (10 años, Entendimiento de 1994);
  \item la UA no eleva en conjunto la protección.
  \end{itemize}
\item \emph{Cláusula de habilitación} (1979): entre países en desarrollo (Mercosur, ASEAN, ALADI). \emph{Art.~V AGCS}: servicios \vertema{3.B.34}.
\item Comité de ACR (1996); \concepto{Mecanismo de Transparencia} (2006). Solo uno declarado compatible (Rep. Checa-Eslovaquia).
\end{itemize}
\graficorepaso[\linewidth]{acr-anuales}{ACR en vigor: 44 (1990), 99 (1995), 390 (2026); máximo de nuevos en 2021 (Reino Unido)}

% ==========================================================================
\bloque[13]{II. Los procesos no comunitarios\label{rep:II}}
\begin{ideaclave}
Integración comercial con poca institución. Funciona donde hay cadenas de valor (T-MEC, ASEAN, RCEP), menos donde se buscaba sustituir importaciones (Mercosur, CAN, África). La ola actual: megaacuerdos que cruzan regiones.
\end{ideaclave}
\graficorepaso[\linewidth]{mapa-bloques}{Los bloques y su superposición: color, bloque regional; trama, CPTPP; contorno, RCEP; número, acuerdos por país}

\sub{II.1. Europa y Eurasia\label{rep:II.1}}
\subsub{a) AELC y EEE}
\begin{itemize}
\item \concepto{AELC} (1960): Islandia, Liechtenstein, Noruega, Suiza; ZLC industrial; se vació hacia la CEE (dominó).
  \begin{itemize}
  \item \concepto{EEE} (1994, sin Suiza): mercado único sin PAC ni UA.
  \item $\approx$30 ALC: India (1-10-2025, 100.000 M\$ de inversión); Mercosur (vigor 1-10-2026 para Islandia y Brasil). Comercio interno 0,4\,\%.
  \end{itemize}
\end{itemize}
\subsub{b) CEFTA, Reino Unido y UEE}
\begin{itemize}
\item \concepto{CEFTA} (1992, rehecho en 2006): Balcanes occidentales y Moldavia; sala de espera de la UE.
\item \emph{Reino Unido}: ACC con la UE (2021); CPTPP (15-12-2024); India (15-7-2026); CCG (firmado 5-2026).
  \begin{itemize}
  \item OBR: comercio $-$15\,\%, productividad $-$4\,\%.
  \end{itemize}
\item \concepto{UEE} (2015): Rusia, Kazajistán, Bielorrusia, Armenia, Kirguistán; UA; más geopolítica que economía.
\end{itemize}

\sub{II.2. América\label{rep:II.2}}
\begin{itemize}
\item \emph{CEPAL}: regionalismo cerrado (sustitución de importaciones) $\to$ abierto (90).
\end{itemize}
\subsub{a) T-MEC (antes TLCAN)}
\begin{itemize}
\item TLCAN 1994; T-MEC 1-7-2020: ZLC profunda.
  \begin{itemize}
  \item Autos: 75\,\% de contenido regional y 40-45\,\% con salarios $\geq$16 \$/h.
  \item Revisión a los 6 años (art.~34.7): 1-7-2026 EE.~UU. no prorroga $\Rightarrow$ revisiones anuales hasta 2036.
  \end{itemize}
\item Aranceles IEEPA 2025 (exentos los que cumplen el origen) anulados (feb. 2026); los de la sec.~232 siguen. México: hasta 50\,\% a países sin acuerdo (2026).
\item 29,5\,\% del PIB mundial; 49\,\% de comercio interno.
\end{itemize}
\subsub{b) Mercosur}
\begin{itemize}
\item Asunción 1991; Ouro Preto 1994: UA imperfecta, AEC $\approx$11\,\%.
  \begin{itemize}
  \item Argentina, Brasil, Paraguay, Uruguay, Bolivia (2024); Venezuela suspendida.
  \end{itemize}
\item UE: firmado 17-1-2026; interino provisional desde 1-5-2026; el de asociación espera al TJUE y a los 27.
\item Debate de \emph{flexibilización}: Uruguay (pide el CPTPP) y Argentina quieren firmar por su cuenta: incompatible con una UA.
\item Comercio interno: 21\,\% (2000) $\to$ 12\,\% (2025): escollo.
\end{itemize}
\subsub{c) Alianza del Pacífico y CAN}
\begin{itemize}
\item \concepto{Alianza del Pacífico} (2011-2012): Chile, Colombia, México, Perú; abierta, profunda, mira a Asia; 2\,\% interno.
\item \concepto{CAN} (Pacto Andino 1969): Bolivia, Colombia, Ecuador, Perú; UA imperfecta; Chile sale 1976, Venezuela 2006; guerra arancelaria Ecuador-Colombia (2026).
\end{itemize}
\subsub{d) Otros}
\begin{itemize}
\item \emph{ALADI} (1980): paraguas de acuerdos parciales. \emph{MCCA}/\emph{SICA}: CAFTA-DR, UE (2012); 27\,\% interno.
\item \emph{CARICOM} (1973): CSME; libre circulación plena de 4 países (1-10-2025).
\item \emph{ALCA} (EE.~UU., desde 1994): muere en Mar del Plata (2005) por Brasil, Argentina y Venezuela; contrapropuesta: \emph{ALBA} (2004).
\end{itemize}
\graficorepaso[0.85\linewidth]{bloques-pib-poblacion}{Peso en el PIB y la población mundiales (2025): los bloques se solapan}

\sub{II.3. Asia-Pacífico\label{rep:II.3}}
\begin{itemize}
\item \emph{Fábrica Asia}: integración de mercado primero, acuerdos después.
\end{itemize}
\subsub{a) ASEAN}
\begin{itemize}
\item 1967; 11 con Timor Oriental (26-10-2025); AFTA (1992), Comunidad (2015), Visión 2045.
  \begin{itemize}
  \item ATIGA mejorado y ACFTA 3.0 (oct. 2025); consenso y no injerencia; 22\,\% interno; centro de los ASEAN+1.
  \end{itemize}
\end{itemize}
\subsub{b) RCEP}
\begin{itemize}
\item Firmado 2020, vigor 2022: ASEAN-10, China, Japón, Corea, Australia, N. Zelanda; la India se fue (2019).
  \begin{itemize}
  \item $\approx$90\,\% de aranceles; \emph{acumulación} de origen; superficial.
  \item Adhesiones: Hong Kong, Sri Lanka, Chile, Bangladés.
  \end{itemize}
\end{itemize}
\subsub{c) CPTPP}
\begin{itemize}
\item TPP (2016) $\to$ EE.~UU. sale (2017) $\to$ CPTPP (30-12-2018): 11 + Reino Unido (2024).
  \begin{itemize}
  \item Costa Rica (cerrada en lo sustancial, mayo 2026); Uruguay; China y Taiwán piden (2021); diálogo con la UE desde nov. 2025.
  \item El más exigente y abierto.
  \end{itemize}
\end{itemize}
\subsub{d) APEC y Asia meridional}
\begin{itemize}
\item \concepto{APEC} (1989): 21 economías, $\approx$60\,\% del PIB; no vinculante; Bogor (1994), Putrajaya 2040.
\item \emph{SAARC/SAFTA}: bloqueada (India-Pakistán); 6\,\% interno. \emph{ANZCERTA} (1983).
\end{itemize}
\begin{center}\small
\newcommand{\si}{\textcolor{tceeverde}{$\checkmark$}}
\newcommand{\no}{\textcolor{tceegris}{--}}
\renewcommand{\arraystretch}{1.1}
\begin{tabular}{@{}>{\raggedright\arraybackslash}p{3.9cm}c@{\ }c@{\ }c@{\ }c@{}}
\textbf{Plato de espaguetis} & \textbf{ASEAN} & \textbf{RCEP} & \textbf{CPTPP} & \textbf{APEC}\\ \hline
Brunéi, Malasia, Singap., Vietnam & \si & \si & \si & \si\\
Indonesia, Filipinas, Tailandia & \si & \si & \no & \si\\
Camboya, Laos, Myanmar & \si & \si & \no & \no\\
Japón, Austral., N. Zelanda & \no & \si & \si & \si\\
China, Corea del Sur & \no & \si & \no & \si\\
Canadá, México, Chile, Perú & \no & \no & \si & \si\\
Reino Unido & \no & \no & \si & \no\\
EE.~UU., Rusia, Taiwán & \no & \no & \no & \si\\ \hline
\end{tabular}
\end{center}

\sub{II.4. Oriente Medio y norte de África\label{rep:II.4}}
\begin{itemize}
\item \concepto{CCG} (1981): 6 monarquías; UA (2003, AEC 5\,\%), mercado común (2008); moneda común parada; Reino Unido (2026); 13\,\% interno.
\item \emph{GAFTA} (1997-2005), 18 países árabes; \emph{UMA} (1989) paralizada; \emph{Agadir} (2004).
\end{itemize}

\sub{II.5. África\label{rep:II.5}}
\subsub{a) Las comunidades regionales}
\begin{itemize}
\item \emph{Obstáculos}: poca complementariedad, infraestructuras, disparidad, conflictos, acuerdos superpuestos. 8 comunidades regionales reconocidas por la UA.
\item \emph{CEDEAO} (1975): salen Mali, Burkina Faso y Níger (2025). \emph{UEMOA} y \emph{CEMAC}: franco CFA \vertema{3.B.16}.
\item \emph{SACU} (1910): la UA más antigua. \emph{EAC}: UA (2005), mercado común (2010), 8 miembros. \emph{Tripartita} COMESA-EAC-SADC (2024).
\end{itemize}
\subsub{b) ZLECAf}
\begin{itemize}
\item Kigali (2018), vigor 2019; 54 firmantes, 50 ratificaciones.
  \begin{itemize}
  \item 90\,\% de líneas en 5 años (10 PMA); Comercio Guiado (2022): 26 países comercian. Interno: 17\,\%.
  \end{itemize}
\end{itemize}

% ==========================================================================
\bloque[6,5]{III. Balance y situación actual\label{rep:III}}
\begin{ideaclave}
Integran mucho donde hay cadenas de valor; en un mundo fragmentado, son la forma de seguir abiertos sin depender de la OMC ni de EE.~UU.
\end{ideaclave}

\sub{III.1. Efectos económicos de los bloques\label{rep:III.1}}
\subsub{a) Medición}
\begin{itemize}
\item \emph{Ex ante}, MEGC; \emph{ex post}, gravedad con resistencia multilateral (\autor{Anderson}-\autor{van Wincoop}, 2003).
\item Un ALC duplica el comercio en 10 años (\autor{Baier}-\autor{Bergstrand}, 2007).
\end{itemize}
\subsub{b) Resultados}
\begin{itemize}
\item \emph{TLCAN} (\autor{Caliendo}-\autor{Parro}, 2015): comercio de México con sus socios +118\,\%.
  \begin{itemize}
  \item Bienestar: México +1,31\,\%, EE.~UU. +0,08\,\%, Canadá $-$0,06\,\%; salarios locales (\autor{Hakobyan}-\autor{McLaren}, 2016).
  \end{itemize}
\item \emph{Mercosur}: desviación (\autor{Yeats}, 1998); peor RRI de los excluidos (\autor{Chang}-\autor{Winters}, 2002).
\item \emph{RCEP} +186.000 M\$ y \emph{CPTPP} +147.000 M\$ de renta mundial en 2030 (\autor{Petri}-\autor{Plummer}, 2020).
  \begin{itemize}
  \item RCEP: +428.000 M\$ de comercio interno frente a $-$48.000 M\$ con terceros: creación $\gg$ desviación.
  \end{itemize}
\item \emph{ZLECAf}: +7\,\% de renta en 2035, 2/3 por facilitación (Banco Mundial, 2020).
\item \emph{Más allá del comercio}:
  \begin{itemize}
  \item IED que produce dentro para todo el bloque (México) \vertema{3.B.18};
  \item convergencia o divergencia (\autor{Venables}, 2003);
  \item \emph{Brexit}, el experimento inverso.
  \end{itemize}
\item \emph{Profundos}: más creación y menos desviación (normas que benefician a terceros; \autor{Mattoo} \emph{et al.}, 2022).
\item \emph{Preferencias sin usar}: solo el 17\,\% del comercio las aprovecha (2022): arancel cero o reglas de origen más caras que el margen. El plato de espaguetis, medido.
\end{itemize}
\subsub{c) Cuánto integran y cuánto pesan}
\begin{itemize}
\item UE (61\,\%) y T-MEC (49\,\%), los más integrados; RCEP y ASEAN, la «fábrica Asia».
  \begin{itemize}
  \item Sur $<$20\,\%; Mercosur, CAN y CARICOM, a la baja desde 2000.
  \item Alianza y AELC: bloques de reglas, no de comercio interno.
  \end{itemize}
\item T-MEC + UE + RCEP $\approx$75\,\% del PIB mundial; ZLECAf y SAFTA: $>$40\,\% de la población con $<$7\,\% del PIB.
\item \emph{Recapitulando}: más creación que desviación donde hay cadenas de valor; ganancias reales, moderadas y desiguales; conviven con la NMF.
\end{itemize}
\graficorepaso[0.85\linewidth]{comercio-intrarregional}{Comercio intrarregional 2000 y 2025: alto en la UE, el T-MEC y el RCEP; bajo y menguante en el Mercosur, la CAN y el CARICOM}

\sub{III.2. Un mundo fragmentado\label{rep:III.2}}
\begin{itemize}
\item \emph{Cambio de contexto}: OMC atascada, sin juez de apelación; EE.~UU. usa aranceles por seguridad, saldo y política exterior.
  \begin{itemize}
  \item NMF: 83\,\% del comercio (2022) $\to$ 72\,\% (2026) \vertema{3.B.33} \vertema{3.B.21}.
  \end{itemize}
\item \emph{4 tendencias}:
  \begin{enumerate}
  \item \emph{EE.~UU., de arquitecto a socio incómodo}: sale del TPP, no renueva el T-MEC; sus pactos arancelarios no son del art.~XXIV. Riesgo \autor{Krugman}: bloque grande con poder de mercado.
  \item \emph{Megaacuerdos como refugio}: CPTPP y RCEP se amplían; la UE firma con Mercosur y México, cierra Indonesia, India, Australia y dialoga con el CPTPP \vertema{3.B.41}: dominó defensivo.
  \item \emph{China como centro}: RCEP, ACFTA 3.0, solicitud al CPTPP; APEC 2026 (Shenzhen) recupera el FTAAP.
  \item \emph{El Sur se integra}: ZLECAf, apertura del Mercosur, adhesiones a ASEAN y RCEP; pero \emph{friend-shoring}: la seguridad pesa más que la eficiencia \vertema{3.B.18} \vertema{3.B.32}.
  \end{enumerate}
\item \emph{¿Peldaño o escollo, hoy?}: no sustituyen a la OMC, pero mantienen vivas las reglas y pueden ser la base de una multilateralización (CPTPP-UE, RCEP-APEC); riesgo: bloques rivales.
\item G-20 y OCDE \vertema{3.B.19}.
\end{itemize}

% ==========================================================================
\bloque[1]{Conclusión}
\begin{itemize}
\item \textbf{Recapitulación}: liberalización discriminatoria.
  \begin{itemize}
  \item Crea comercio entre socios y puede desviarlo; la OMC la admite con condiciones poco exigentes.
  \item \autor{Krugman}: pocos bloques grandes pueden ser lo peor; \autor{Baldwin}: contagio, peldaño si abierto; \autor{Krishna}: las rentas bloquean.
  \item Fuera de la UE, ZLC con poca institución; funcionan donde hay cadenas de valor (T-MEC, ASEAN, RCEP), no donde se sustituían importaciones.
  \item Más creación que desviación; ganancias moderadas y desiguales.
  \end{itemize}
\item \textbf{Relevancia}: con la OMC atascada, la vía de la liberalización; para España, diversificar \vertema{4.A.24}.
\item \textbf{Extensiones}: teoría de la integración \vertema{3.B.10}; OMC y art.~XXIV \vertema{3.B.33}; UE \vertema{3.B.36} y sus acuerdos \vertema{3.B.41}; áreas emergentes \vertema{3.B.32}; empresas españolas en América Latina y Asia \vertema{4.A.23}.
\item \textbf{Opinión}: el riesgo no es el plato de espaguetis, sino el poder de mercado de los grandes.
  \begin{itemize}
  \item Mejor seguro: bloques abiertos y profundos (CPTPP); los cerrados (Mercosur) perdieron la apuesta.
  \end{itemize}
\item \textbf{Idea final}: ¿más vale solo? Casi nunca.
  \begin{itemize}
  \item Las cadenas de valor han sido buena compañía.
  \item Abiertos (nuevos socios, aranceles bajos): peldaño; cerrados: los 3 bloques que temía \autor{Krugman}.
  \end{itemize}
\end{itemize}

% ==========================================================================
\anexorepaso
\bloque{Datos clave}
\begin{itemize}
\item ACR en vigor: \dato{390}{OMC, nov.\ 2026}; $\geq$84 sin notificar.
\item PIB mundial: \dato{T-MEC 29,5\,\%; RCEP 27,2\,\%; UE 17,9\,\%; CPTPP 14,3\,\%}{2025}.
\item Comercio interno: \dato{UE 61\,\%, T-MEC 49\,\%, RCEP 38\,\%, Mercosur 12\,\%}{2025}.
\item NMF: \dato{83\,\% (2022) $\to$ 72\,\% (2026)}{OMC}.
\end{itemize}
\bloque{Autores y fechas}
\hito{1950}{\autor{Viner}: creación y desviación}
\hito{1961}{\autor{Balassa}: formas de integración}
\hito{1976}{\autor{Kemp}-\autor{Wan}}
\hito{1979}{Cláusula de habilitación}
\hito{1991}{\autor{Krugman}: 3 bloques}
\hito{1993}{\autor{Baldwin}: dominó}
\hito{1995}{\autor{Bhagwati}: plato de espaguetis}
\hito{1998}{\autor{Krishna}: intereses creados}
\hito{2015}{\autor{Caliendo}-\autor{Parro}: TLCAN}
\bloque{Preguntas del tribunal}
\begin{enumerate}
\item \textbf{¿ZLC o UA?} La ZLC necesita reglas de origen; la UA, política comercial común.
\item \textbf{¿Por qué 3 bloques es lo peor?} Más comercio libre frente a más poder de mercado (\autor{Krugman}).
\item \textbf{¿Por qué un ACR puede frenar la OMC?} Crea rentas de preferencia que se defienden (\autor{Krishna}; \autor{Limão}).
\item \textbf{¿Qué exige el art.~XXIV?} Lo esencial del comercio, plazo razonable, no subir la protección.
\item \textbf{¿RCEP o CPTPP?} RCEP: grande y superficial, con China; CPTPP: profundo y abierto, sin China ni EE.~UU.
\item \textbf{¿Qué pasa con el T-MEC tras julio de 2026?} Revisiones anuales; expira en 2036 salvo renovación.
\end{enumerate}
\begin{trampa}
\begin{itemize}
\item AELC: \emph{4} miembros; Suiza \emph{no} está en el EEE.
\item El RCEP \emph{no} incluye a la India; el CPTPP \emph{no} incluye a EE.~UU. ni a China.
\item ASEAN: \emph{11} miembros (Timor Oriental, 2025).
\item Mercosur, ASEAN y ALADI: \emph{habilitación}, no art.~XXIV.
\end{itemize}
\end{trampa}

\end{columnas}
\end{document}
